\section{Research Experience}

\begin{research}
{Starspot Inference with Light Curve Inversion Techniques}
{American Museum of Natural History}
{Dr. Lucy Lu}
{Aug 2023 -- Sep 2025}
    \item Investigated the feasibility and limitations of reconstructing stellar surface features from rotational photometric variability.
    \item Applied light-curve inversion, time-series analysis, and Bayesian inference to study non-unique starspot configurations and model degeneracies.
    \item Developed Python-based computational workflows for simulation, model fitting, and evaluation as part of an M.S. thesis on starspot inference.
\end{research}

\begin{research}
{Towards Mapping Brown Dwarf and Giant Exoplanet Atmospheres}
{American Museum of Natural History}
{Dr. Johanna Vos}
{May 2021 -- Aug 2023}
    \item Investigated atmospheric variability in brown dwarfs and giant exoplanets using rotational photometric observations.
    \item Applied time-series analysis and statistical modeling to characterize cloud-driven variability and infer atmospheric structure.
    \item Developed Python workflows for processing observational datasets and analyzing photometric light curves.
\end{research}

\begin{research}
{Chemodynamically Characterizing the Jhelum Stellar Stream with APOGEE-2}
{Sloan Digital Sky Survey Faculty and Student Team}
{Dr. Allyson Sheffield}
{May 2020 -- May 2021}
    \item Investigated the chemical and dynamical properties of candidate stellar members of the Jhelum stellar stream using APOGEE-2 spectroscopy.
    \item Applied statistical analysis to identify stream members and characterize their chemical abundance patterns.
    \item Contributed to a refereed publication in \textit{The Astrophysical Journal}.
\end{research}

\begin{research}
{Looking at Star Formation Through Chemistry}
{American Museum of Natural History}
{Dr. Allyson Sheffield}
{Aug 2018 -- May 2020}
    \item Investigated stellar populations through spectroscopic measurements of chemical abundances.
    \item Processed and normalized echelle spectra and measured equivalent widths of spectral lines.
    \item Applied quantitative spectroscopic analysis techniques to characterize stellar chemical compositions.
\end{research}

\begin{research}
{Two-Point Statistics in the Star Forming ISM}
{Flatiron Institute, Center for Computational Astrophysics}
{Dr. Chang-Goo Kim}
{Jun 2018 -- Aug 2018}
    \item Investigated statistical relationships between metallicity, gas dynamics, and star formation in the interstellar medium.
    \item Applied two-point statistical methods to analyze correlations within numerical simulation data.
\end{research}